%=====================================================================
\chapter{Why Not Just Files?}\label{ch:whynotfiles}
%=====================================================================
\locator{1}{Part I --- What a Database Is}

\begin{outcomes}
By the end of this chapter you will be able to:
\begin{tight}
  \item \textbf{Describe} how a file-based application stores and finds data,
        and \textbf{state} what it costs to add one field.
  \item \textbf{Name} the six problems a file-based approach has, and
        \textbf{give} a concrete example of each.
  \item \textbf{Explain} what a database and a database management system are,
        and \textbf{distinguish} the two.
  \item \textbf{Judge} when a file is the right answer, which is more often
        than this chapter's title suggests.
  \item \textbf{Predict} which of the six problems a given scenario will hit
        first.
\end{tight}
\end{outcomes}

\begin{prereq}
One semester of programming: reading and writing a file, loops, and what a
record looks like when it is a line of comma-separated text. Nothing else.
Appendix~A if you have not checked that SQLite is on your machine.
\end{prereq}

\begin{keyterms}
data $\bullet$ information $\bullet$ file-based system $\bullet$ database
$\bullet$ database management system (DBMS) $\bullet$ data redundancy
$\bullet$ data inconsistency $\bullet$ data isolation $\bullet$ integrity
constraint $\bullet$ atomicity $\bullet$ concurrent access $\bullet$ lost
update $\bullet$ data security $\bullet$ application program
\end{keyterms}

\section{The Program You Have Already Written}

Somewhere in your first semester you wrote a program that kept a list of
things in a file. It probably looked like this, and there is nothing wrong
with it:

\begin{lstlisting}[language=python]
# students.csv
# 2021-IT-001,Ayesha Khan,BSIT,3
# 2021-IT-002,Bilal Ahmed,BSIT,3

def find(roll_no):
    for line in open("students.csv"):
        fields = line.strip().split(",")
        if fields[0] == roll_no:
            return fields
    return None
\end{lstlisting}

It works. It is twelve lines. For a list of your own books, or a configuration
file, or the output of a program that nobody else reads, it is the right
answer and this chapter is not an argument against it.

This chapter is about the afternoon it stops being the right answer --- and
about the fact that it does not stop gradually. A file-based approach works,
and works, and works, and then the department asks for one more thing and
suddenly nothing works.

\begin{definitionbox}[Data, Information, and the difference]
\term{Data} are raw recorded facts: the characters \texttt{2021-IT-001} in a
file.

\smallskip
\term{Information} is data in a context that makes it answer a question:
\emph{the roll number of the student who scored highest in Database Systems
last term}.

\smallskip
A \term{file-based system} stores data and leaves every question to the
application program. A \term{database} stores data \emph{together with a
description of what it means and what must be true of it}, so that some of the
questions --- and all of the guarantees --- can be answered by the store
itself.

\smallskip
That last sentence is the whole subject. Everything in the remaining
twenty-two chapters is a particular case of it.
\end{definitionbox}

\section{Six Things That Go Wrong}

These are not six opinions. Each is a specific failure with a name, and you
will meet all six in a working life.

\dsfig{%
\begin{tikzpicture}[
  app/.style={rounded corners=3pt, draw=#1, line width=1.0pt, fill=#1!8,
              text=#1!60!black, font=\bfseries\scriptsize, align=center,
              text width=2.5cm, minimum height=0.9cm},
  fil/.style={draw=neutral, line width=0.9pt, fill=neutral!8,
              text=neutral, font=\scriptsize, align=center,
              text width=2.3cm, minimum height=0.8cm},
  arr/.style={-{Stealth[length=2.0mm]}, neutral!70, line width=0.8pt},
  note/.style={font=\scriptsize\itshape, text=accent, align=center,
               text width=5.0cm}
]
  % --- file-based, on the left
  \node[font=\bfseries\small, text=accent] at (3.0,3.5) {File-based};
  \node[app=accent] (a1) at (0.6,2.5) {Enrolment\\program};
  \node[app=accent] (a2) at (3.0,2.5) {Results\\program};
  \node[app=accent] (a3) at (5.4,2.5) {Fees\\program};
  \node[fil] (f1) at (0.6,0.9) {students.csv};
  \node[fil] (f2) at (3.0,0.9) {marks.csv};
  \node[fil] (f3) at (5.4,0.9) {fees.csv};
  \draw[arr] (a1) -- (f1);
  \draw[arr] (a2) -- (f2);
  \draw[arr] (a3) -- (f3);
  \node[note] at (3.0,-0.3)
       {Each program owns its own file and its own idea of what a student is.
        The name is stored three times.};

  % --- database, on the right
  \node[font=\bfseries\small, text=greenhl] at (11.4,3.5) {Database};
  \node[app=greenhl] (b1) at (9.0,2.5) {Enrolment\\program};
  \node[app=greenhl] (b2) at (11.4,2.5) {Results\\program};
  \node[app=greenhl] (b3) at (13.8,2.5) {Fees\\program};
  \node[rounded corners=3pt, draw=greenhl, line width=1.2pt, fill=greenhl!14,
        text=greenhl!60!black, font=\bfseries\scriptsize, align=center,
        text width=5.6cm, minimum height=0.9cm] (dbms) at (11.4,1.3) {DBMS};
  \node[fil, text width=5.6cm, draw=greenhl] (db) at (11.4,0.1)
       {one database};
  \foreach \b in {b1,b2,b3}{\draw[arr] (\b) -- (dbms);}
  \draw[arr] (dbms) -- (db);
  \node[note, text=greenhl!60!black] at (11.4,-1.1)
       {One store, one description of what a student is, and one place the
        rules are kept.};
\end{tikzpicture}}%
{The same three programs, two ways. The difference is not that the right-hand
side has fewer files --- it is that the rules about what a student is have
moved out of the programs and into one place.}{filesvsdb}

\begin{alertbox}[The six, with a case for each]
\begin{tightnum}
  \item \textbf{Data redundancy.} The student's name is in
        \texttt{students.csv}, in \texttt{marks.csv} so the results sheet can
        print it, and in \texttt{fees.csv} so the receipt can. Three copies,
        and disk space is the least of it.
  \item \textbf{Data inconsistency.} A student marries and changes her name.
        Two of the three files are updated. The college now holds two different
        answers to one question, and \emph{nothing in the system is wrong} ---
        each file is internally fine. This is redundancy's actual cost.
  \item \textbf{Data isolation.} The question ``which students owe fees and
        are enrolled in a course that has already started?'' needs two files in
        different formats. There is no way to ask it; there is only a program
        somebody has to write.
  \item \textbf{Integrity problems.} ``A student's semester is between 1 and
        8.'' Where does that rule live? In the enrolment program. And in the
        results program, where somebody wrote $1$ to $10$. The rule is a
        comment in three places and enforced in none.
  \item \textbf{Atomicity problems.} Transferring a student between sections
        means removing a line from one file and adding it to another. The
        machine loses power in between. The student is now in no section at
        all, and nothing can tell you that happened.
  \item \textbf{Concurrent access and security.} Two clerks open the file, both
        edit, both save; the second overwrites the first and neither is told
        --- a \term{lost update}. And the fees clerk, who needs one column, has
        the whole file including every home address.
\end{tightnum}
\end{alertbox}

\begin{conceptbox}[Which one bites first, and why it is never the one you
expect]
Students asked to rank these almost always put redundancy first, because it is
the one that is visible --- you can \emph{see} the duplicated column.

\smallskip
In practice the order is almost exactly reversed. Redundancy costs disk, which
is cheap. \textbf{Inconsistency} is what gets noticed, because somebody
receives a letter addressed to a name they no longer use. \textbf{Isolation}
is what gets the department to pay for a new system, because the question they
want answered cannot be answered at all. And \textbf{atomicity} is the one
nobody finds until the day it matters, because a half-finished transfer looks
exactly like a student who was never enrolled.

\smallskip
Keep this ranking in mind in \chref{project}, when you have to justify your
design to somebody who will ask what problem it solves.
\end{conceptbox}

\section{What a Database Management System Is}

\begin{definitionbox}[Database, DBMS, Database System]
A \term{database} is a shared, organised collection of data, together with a
description of that data --- what it is called, what type it is, what must be
true of it, and how its parts relate.

\smallskip
A \term{database management system} (DBMS) is the \emph{software} that creates,
maintains and controls access to a database. SQLite is a DBMS; so are Oracle,
PostgreSQL, MySQL and SQL Server.

\smallskip
A \term{database system} is the whole arrangement: the data, the DBMS, the
application programs that use it, and the people.

\smallskip
These three are routinely used as if they meant the same thing, including by
people who should know better. In an examination they do not, and the
distinction is CLO-1's business.
\end{definitionbox}

The DBMS is what the right-hand side of Figure~\ref{fig:filesvsdb} adds. It is
a single piece of software that every program goes through, and because
everything goes through it, it is the one place where a rule can be stated once
and be true for everybody.

\begin{examplebox}[What moves into the DBMS, problem by problem]
\begin{tabular}{@{}L{4.6cm}L{9.6cm}@{}}
\textbf{Redundancy} & One copy of each fact, and other tables refer to it
(\chref{mapping}) \\
\textbf{Inconsistency} & Follows from the above --- there is nothing to
disagree with \\
\textbf{Isolation} & One query language over all the data (\chref{algebra}
onward) \\
\textbf{Integrity} & Constraints declared once, in the schema
(\chref{constraints}) \\
\textbf{Atomicity} & Transactions: all of it or none of it
(\chref{transactions}) \\
\textbf{Concurrency, security} & Locking, and authorisation
(\chref{transactions}, \chref{project}) \\
\end{tabular}

\smallskip
That table is this book's table of contents, rearranged. Every chapter from
here is one row of it, done properly.
\end{examplebox}

\section{When a File Is Still the Right Answer}

A course called Database Systems has an obvious bias, and you should know
about it.

\begin{conceptbox}[Use a file when every one of these is true]
\begin{tight}
  \item \textbf{One program reads it, and one writes it.} No concurrency, so
        no lost updates.
  \item \textbf{It is small enough to read entirely.} If you always want all of
        it, an index buys nothing.
  \item \textbf{The questions never change.} Isolation only hurts when somebody
        asks something new.
  \item \textbf{Losing it is survivable}, or it can be regenerated.
\end{tight}

\smallskip
Configuration files, log files, a cache, the output of a nightly job --- all
files, all correctly so. \textbf{The failure is not using a file; it is using a
file after the fourth condition stopped holding and not noticing.}
\end{conceptbox}

\begin{pitfall}
\begin{tight}
  \item \textbf{Treating redundancy as the main problem.} It is the cheapest of
        the six and the easiest to see, which is why it gets the attention.
  \item \textbf{Saying ``database'' when you mean ``DBMS''.} The data is not
        the software.
  \item \textbf{Believing a DBMS removes the need to think about integrity.}
        It gives you somewhere to put the rules. It does not write them, and
        \chref{constraints} measures what happens when you declare one and it
        is not enforced.
  \item \textbf{Assuming a database is always better.} Four conditions above
        say otherwise.
  \item \textbf{Calling a spreadsheet a database.} It has rows and columns and
        none of the six guarantees.
\end{tight}
\end{pitfall}

\begin{lab}[Break a file-based system on purpose]
No database yet. This laboratory is about earning the next twenty-two
chapters.

\begin{lstlisting}[language=bash]
# --- 1. Make the files -----------------------------------------------
# Three CSV files, by hand, in a text editor. Five students each.
#   students.csv   roll_no,name,programme,semester
#   marks.csv      roll_no,name,course,marks
#   fees.csv       roll_no,name,amount,paid

# --- 2. Count the redundancy -----------------------------------------
# How many times does one student's NAME appear across the three files?
# Multiply by 4,000 students. Write the number down.

# --- 3. Create an inconsistency --------------------------------------
# Change one student's name in students.csv and marks.csv but NOT in
# fees.csv. Now answer: what is that student's name? Justify it.

# --- 4. Hit data isolation -------------------------------------------
# Write the program that answers: "which students owe fees AND are
# enrolled in a course where they scored below 50?"
# Time yourself. Do not finish it if it takes more than twenty minutes
# -- stopping is the result.

# --- 5. Break atomicity ----------------------------------------------
# Write a script that moves a student from one file to another, with a
# deliberate crash between the two writes. Run it. Describe the state.

# --- 6. Simulate a lost update ---------------------------------------
# Open students.csv in two editors. Change a different row in each.
# Save both. Report what survived.

# --- 7. Write the memo ------------------------------------------------
# One page to a head of department who has asked why the existing
# system should be replaced. Use YOUR numbers from steps 2 and 4.
# No jargon: the reader has not taken this course.
\end{lstlisting}

\smallskip
\textbf{What to notice.} Step 4 is the one that convinces people. Redundancy is
an argument; a question that cannot be answered at all is a business problem,
and it is the one that gets a project funded.
\end{lab}

\begin{checkpoint}
\begin{tightnum}
  \item A college keeps student records in one file and library loans in
        another, each maintained by one clerk who never touches the other.
        Which of the six problems does this arrangement have, and which does it
        genuinely escape? Say what would have to change for your answer to
        change.
  \item Distinguish a database from a database management system, and explain
        why ``we bought a database'' is a sentence that needs clarifying before
        you can act on it.
  \item Your colleague argues that redundancy is harmless because disk is
        cheap. Give the argument you would make in reply, name the problem
        redundancy actually causes, and give an example in which no single file
        is wrong and the system still is.
\end{tightnum}
\end{checkpoint}

\begin{chaptersummary}
\begin{tight}
  \item Data are recorded facts; information is data in a context that answers
        a question. A database stores data \emph{with a description of what it
        means and what must be true of it}, which is the whole subject.
  \item A file-based system works until it does not, and the change is abrupt
        rather than gradual.
  \item The six problems are redundancy, inconsistency, isolation, integrity,
        atomicity, and concurrent access with security.
  \item Redundancy is the visible one and the cheapest. Inconsistency is what
        gets noticed, isolation is what gets a project funded, and atomicity is
        the one nobody finds until it matters.
  \item A \emph{database} is the data; a \emph{DBMS} is the software; a
        \emph{database system} is the whole arrangement including the people.
  \item Everything goes through the DBMS, which is why it is the one place a
        rule can be stated once and be true for everybody.
  \item A file is still right when there is one reader, one writer, a fixed set
        of questions and nothing irreplaceable. The failure is not noticing
        when that stops being true.
  \item The table in Section 1.3 is this book's contents rearranged: each of
        the six problems is solved by a named later chapter.
\end{tight}
\end{chaptersummary}

\begin{reviewq}
  \item \textbf{(MCQ)} Two files hold a student's name and only one is updated.
        This is an example of: (a)~redundancy; (b)~inconsistency; (c)~isolation;
        (d)~an integrity constraint.
  \item \textbf{(MCQ)} The software that creates and controls access to a
        database is called: (a)~the database; (b)~the database system; (c)~the
        DBMS; (d)~the schema.
  \item \textbf{(Short)} Give a question about a university that cannot be
        answered from two separate CSV files, and say which of the six problems
        that is.
  \item \textbf{(Short)} Explain atomicity using an example that is not a bank
        transfer.
  \item \textbf{(Short)} State two conditions under which a plain file is a
        better choice than a database, and give a real example of each.
  \item \textbf{(Applied)} A clinic keeps appointments in a spreadsheet shared
        on a network drive, edited by three receptionists. Identify every one of
        the six problems it has, rank them by which will cause a complaint
        first, and justify your ranking.
  \item \textbf{(Applied)} You are asked to justify replacing a working
        file-based system. Write the three strongest arguments in order, say
        what evidence you would gather for each before the meeting, and state
        which argument you would drop if you only had five minutes.
\end{reviewq}
